\section{Social Impact \& Conclusion}
\label{sec:conclusion}

In this work, we contribute a family of models (i.e., \mfourttwo, \expressive, and \streaming) designed to provide end-to-end multilingual, expressive, and streaming translations, marking a pivotal step towards naturalistic S2ST. First, \mfourttwo, an improved version of \mfourt, is the foundational multilingual and multimodal model on which the latter two models are initialized. \expressive enables translation that preserves prosody and the style of one's voice, and \streaming leverages the Efficient Monotonic Multihead Attention (EMMA) mechanism to generate low-latency target translations without waiting for complete source utterances. To evaluate these models, we combined novel and modified versions of existing metrics to evaluate performance and robustness. Taking a four-pronged approach to ensure the safety of our systems, we implemented the first known red-teaming effort for multimodal machine translation, an added-toxicity detection and mitigation system, evaluations for gender bias, and an inaudible localized watermarking mechanism designed to dampen the impact of deepfakes. Consequently, we bring \expressive, and \streaming together as \seamless, the first publicly available system that unlocks expressive cross-lingual communication in real time.

\subsection{Naturalistic Translations \& Experiential Futures}

The downstream applications our family of models may give rise to could meaningfully transform how cross-lingual communication and experiences manifest across online and offline contexts. While recognizing that these technical building blocks could be combined to enable a smorgasbord of experiences, we briefly explore some potential recipes below.

To start, here are a few ways to use our unified system—\seamless:
\begin{itemize}
  \item \textit{Audio or video-calling.} Integrating our models into computer-mediated communication applications or other voice or video-calling platforms would enable real-time and expressive multilingual dialogue.
  For such an experience, a user could select their desired output language, and every utterance in that call made by any speaker would be translated into the pre-determined target language. With expressive translations, listeners should not have trouble differentiating who said what. Moreover, because \mfourttwo supports both speech and text modalities, users should also be able to see live captions alongside speech outputs.
  \item \textit{Augmented or virtual reality (AR/VR) environments}. Akin to the use case above, \seamless could support multi-person cross-lingual interactions in AR/VR environments, be it multiplayer games or conference meetings.
  \item \textit{Online streaming}. \seamless can also be adapted to run locally on users' personal computers. Through a simple interface, one could use the model to listen to and render translated audio outputs to a virtual stream. This would allow users to speak in any language on their streaming platforms.
  \item \textit{Wearable devices: earbuds or smart glasses}. To create the first generation of a real-world Universal Speech Translator, \seamless could be integrated into earbuds, smart glasses, or comparable wearable devices of a similar nature. Today, most earbuds and smart glasses come with compact microphone and speaker systems, which could easily facilitate receiving and producing speech input and output. If every interlocutor in a conversation is equipped with a \seamless-supported device, all parties can speak in whatever language they want and remain comprehensible. Additionally, live captions (as supported by \mfourttwo) could also be displayed on lenses of smart glasses for boosted accessibility and user confidence.
  

  \end{itemize}


Beyond these synchronous use cases, our models could also be used for making passive content consumption more inclusive:
\begin{itemize}
  \item \textit{Translations for voice-messaging platforms.} Today, instead of exclusively relying on texts, many people record audio notes on computer-mediated communication platforms like WhatsApp or Messenger to get their messages across. \mfourttwo and \expressive can support not just the translation of these audio notes, they can do so while preserving prosody and the style of one's voice.
  \item \textit{Long-form audio translation pipeline.} Our models could be incorporated into a larger audio processing pipeline to translate long-form audio content such as lectures or podcasts. By first isolating the voice audio (i.e., removing all music, background noise, and sound effects), the resulting clips can be processed independently by our models before being brought back together to form a fully expressive and multilingual product.
  \item \textit{Video dubbing translation pipeline}. Building on the abovementioned pipeline, combining our models with a video dubbing tool [e.g., Wav2Lip \citep{prajwal2020lip}] could streamline the creation of multilingual video content that is of high quality both on the visual and audible fronts.
\end{itemize}

Overall, the multidimensional experiences \seamless may engender could lead to a step change in how machine-assisted cross-lingual communication is accomplished. For the immigrant interviewees featured in \Cref{sec:problem}, the communicative capabilities \seamless affords may unlock new possibilities in their personal and professional lives. In the near future, with the help of \seamless, everyday communication that once appeared challenging may become ordinary. While not a panacea for social integration, giving them a tool that softens the effects of language barriers could streamline their day-to-day lives in their receiving society and allow them to better pursue personal goals. By publicly releasing our work, we hope that researchers and developers can expand the impact of our contributions by building technologies aimed at bridging multilingual connections in an increasingly interconnected and interdependent world.

\subsection{Ethical Considerations \& Future Work}

Although we built the family of \seamless models to be used widely, we recognize that user populations are heterogeneous and that our systems may work better for some over others \citep{wang2023human}. Despite carefully evaluating our artifacts across multiple fairness axes and implementing checks and balances whenever possible, model performance may vary depending on users' race, accent, or gender. Moreover, because of the dependencies involved, performance gaps at the ASR stage [which have been well documented \citep{koenecke2020racial,ngueajio2022hey}], may lead to subsequent performance degradation at the expressive and streaming levels. As such, some users may have to continue altering their regular speech patterns to take full advantage of the capabilities offered by our models. 

Moreover, as with all technologies, our models are not impervious to unintended use. In the context of \seamless, bad actors could use our models to enact voice phishing (i.e., pretending to be someone on the phone to exploit unsuspecting individuals for money or personal information) or deepfakes. By instigating a watermarking mechanism and releasing its detector, we offer one solution to help users identify the synthetic origin of the content they are potentially exposed to. That said, taming the effects of the malicious use of AI systems requires a multifaceted approach. Alongside individual-level AI literacy and scam prevention tactics, we believe that increasing public awareness and industry-wide standards around such issues are imperative for the safe implementation of comparable systems in the future.

To further the goal of realizing the Universal Speech Translator, future research should continue focusing on improving language coverage and closing the performance gaps between high-resource and low-resource languages. More resources should also be directed at ensuring that emerging systems work well for diverse user groups, especially those that have been historically underprioritized when it comes to AI development. Beyond spoken and written languages, the visual modality in cross-lingual communication, which comprises sign languages and other visual signals (i.e., gestures, facial expressions, lip movements, etc.), deserves further attention. For one, access to all modalities is not always possible—availability may be limited either due to physical (i.e., loss of hearing due to old age) or circumstantial reasons (i.e., reduced capacity for speech at a loud bar). As such, developing integrative and multimodal translation technologies may propel research on adaptive systems that enhance certain modalities when others are compromised. The ability to complement human communication in such a manner not only makes translation tools more robust, it paves the way for a more inclusive and accessible technological future for many more people.